\scanpage{19}
\mspage{- IV-\handwritten{953} -}
En effet , $Y$ étant compact \deleted{\illegible{1}} et noethérier , toute partie de $Y$ est compacte\handwritten{,}
donc fermée , ce qui implique que $Y$ est discret , donc fini puisqu'il est
compact .

\uline{Co\deleted{\unsure{b|o}}rollaire} (18.6.10) .- \uline{Soit} $A$ \uline{un anneau local noethérien . Pour que} $ {}^{\mathrm h}A$ \uline{poss\handwritten{è}de une des propriétés suivantes} :

a) \uline{être un anneau de Cohen-Macaulay} ;

b) \uline{vérifier la propriété} $(S_n)$ ;
\margin{\handwritten{\unsure{virg.|ví}}}

c) \uline{être régulier} ;

d) \uline{vérifier la propriété} $(R_n)$ ;

\uline{il faut et il suffit que} $A$ \uline{possède cette même propriété} .

Cela résulte aussitôt de ce que les fibres du morphisme $\operatorname{Spec}({}^{\mathrm h}A)\to\operatorname{Spec}(A)$
sont géométriquement régulières (18.6.\unsure{9|8}) et de (6.4.1),(6.5.1) et (6.5.3).

\uline{Corollaire} (18.6.11) .- \uline{Soit} $A$ \uline{un anneau semi-local noethérien . Pour qu'un
idéal premier} $\underline p$ \uline{de} $ {}^{\mathrm h}A$ \uline{appartienne à} $\operatorname{Ass}({}^{\mathrm h}A)$ , \uline{il faut et il suffit que} $\underline p\cap A$
\uline{appartienne à} $\operatorname{Ass}(A)$ .

Compte tenu de la description des fibres du morphisme $\operatorname{Spec}({}^{\mathrm h}A)\to\operatorname{Spec}(A)$ ,
cela résulte aussitôt de (3.3.1) .

\uline{Proposition} (18.6.\unsure{12}) .- \uline{Soi\deleted{en}t} $A$ \uline{un anneau local ; les propriétés suivantes
sont équivalentes} :

a) $A$ \uline{est unibranche}\margin{\handwritten{\unsure{,,|virg.}}} (resp. \uline{réduit et unibranche}) .

b) \uline{Pour toute} $A$-\uline{algèbre locale strictement essentiellement étale} $A'$ ,\deleted{\illegible{1}}
$\operatorname{Spec}(A')$ \uline{est irréductible} (resp. \uline{intègre}) .

c) $\operatorname{Spec}({}^{\mathrm h}A)$ \uline{est irréductible} (resp.\uline{intègre}) .

Notons d'abord que $ {}^{\mathrm h}(A_{\mathrm{red}})=({}^{\mathrm h}A)_{\mathrm{red}}$ , car \deleted{\illegible{4}
\unsure{$A_{\mathrm{red}}$} \illegible{4}} $ {}^{\mathrm h}(A_{\mathrm{red}})=({}^{\mathrm h}A)\otimes_A A_{\mathrm{red}}$ (18.6.8) , et
comme $ {}^{\mathrm h}(A_{\mathrm{red}})$ est réduit (18.6.9) , il est égal à $({}^{\mathrm h}A)_{\mathrm{red}}$ . Ceci permet ,
dans la démonstration , de considérer seulement le cas où $A$ est \uline{réduit} ,donc
aussi $A'$ (17.5.7) et $ {}^{\mathrm h}A$ (18.6.9) .

\deleted{\illegible{2}} Avec les notations de (18.6.4) , la condition b) signifie
que tous les $A_{\handwritten{\alpha}}$ sont intègres ; \deleted{\unsure{comme les homomorphismes} \unsure{$A_{\alpha}$ à $A_{\beta}$} sont
\unsure{injectifs} ,} on en conclut que $ {}^{\mathrm h}A=\varinjlim A_{\handwritten{\alpha}}$ est intègre (5.13.3 ) , donc b) entraîne
c) .
% archive sheet number: 18
